% Generated by replication/build_counterfactual_tables.py.
\begin{table}[p]
  \centering
  \caption{Income scenarios with alternative baselines and destinations}
  \label{tab:income-counterfactual-sensitivity}
  \small
  \begin{tabular}{>{\raggedright\arraybackslash}p{0.52\textwidth}cc}
    \toprule
    Scenario & Predicted baseline & Observed baseline \\
    & (\%) & (\%) \\
    \midrule
    Informal to formal, same size & 13.5 & 13.3 \\[0.4em]
    Formal workers to formal 101+ & 3.8 & 3.8 \\[0.4em]
    Informal workers to informal 101+ & 7.6 & 7.5 \\[0.4em]
    All workers to formal 101+ & 25.9 & 25.5 \\[0.4em]
    \midrule
    \multicolumn{3}{l}{\textit{Destinations selected from the estimated profiles}} \\
    Informal workers to informal 51--100 & 10.1 & 9.9 \\[0.4em]
    Formal to 101+, informal to 51--100 & 13.8 & 13.7 \\[0.4em]
    \bottomrule
  \end{tabular}
  \vspace{0.4em}
  \begin{minipage}{0.97\textwidth}
  \footnotesize
  Notes: Percentage changes in total weekly labor income for the entire 2025 sample. The predicted baseline uses Equation~(\ref{eq:income-counterfactual}). The observed baseline rescales observed income by the same estimated wage ratios, preserving log residuals. Both use pooled coefficients from Equation~(\ref{eq:main}). In the last two rows, destinations maximize the estimated size component within each formality group, including moves from informal 101+ to informal 51--100.
  \end{minipage}
\end{table}
